% Author's solution, images/HW04/IMG_9176.HEIC--IMG_9182.HEIC.
% Use the final parametrization in IMG_9177 and the completed image proof
% in IMG_9180--IMG_9181; omit crossed-out calculations.
\begin{proof}[Solution]
\begin{enumerate}[(a)][2]
  \item The kernel consists of the solutions of \(Ax=0\). Subtracting the
    first row from the third, and then the second row from the third, gives
    \[
      A\sim
      \begin{pmatrix}
        1&2&0&1\\
        0&1&1&1\\
        0&0&0&0
      \end{pmatrix}.
    \]
    Thus the equations are
    \[
      x_1+2x_2+x_4=0,
      \qquad
      x_2+x_3+x_4=0.
    \]
    Taking \(x_3\) and \(x_4\) as free variables, we obtain
    \(x_2=-x_3-x_4\) and \(x_1=2x_3+x_4\). Consequently,
    \[
      x=
      x_3\begin{pmatrix}2\\-1\\1\\0\end{pmatrix}
      +x_4\begin{pmatrix}1\\-1\\0\\1\end{pmatrix},
    \]
    and
    \[
      \ker T=
      \operatorname{span}\left\{
        \begin{pmatrix}2\\-1\\1\\0\end{pmatrix},
        \begin{pmatrix}1\\-1\\0\\1\end{pmatrix}
      \right\}.
    \]
    These two vectors span the kernel. They are linearly independent,
    since a linear combination equal to zero has both coefficients
    zero. Hence they form a basis of \(\ker T\).

  \item Let \(c_1,c_2,c_3,c_4\) denote the columns of \(A\). They satisfy
    \[
      c_2=2c_1+c_3,
      \qquad
      c_4=c_2-c_1=c_1+c_3.
    \]
    For every \(x\in\mathbb{R}^4\),
    \begin{align*}
      Ax
        &=x_1c_1+x_2c_2+x_3c_3+x_4c_4\\
        &=(x_1+2x_2+x_4)c_1+(x_2+x_3+x_4)c_3.
    \end{align*}
    Therefore \(\operatorname{im}T\subseteq
    \operatorname{span}\{c_1,c_3\}\).

    Conversely, if \(y=ac_1+bc_3\), choose
    \(x=(a,0,b,0)^T\). Then \(Ax=y\), so
    \(\operatorname{span}\{c_1,c_3\}\subseteq\operatorname{im}T\).
    We conclude that
    \[
      \operatorname{im}T=
      \operatorname{span}\left\{
        \begin{pmatrix}1\\0\\1\end{pmatrix},
        \begin{pmatrix}0\\1\\1\end{pmatrix}
      \right\}.
    \]
    To check linear independence, suppose \(ac_1+bc_3=0\). Then
    \[
      a\begin{pmatrix}1\\0\\1\end{pmatrix}
      +b\begin{pmatrix}0\\1\\1\end{pmatrix}
      =\begin{pmatrix}a\\b\\a+b\end{pmatrix}=0,
    \]
    which gives \(a=b=0\). Thus \((c_1,c_3)\) is a basis of
    \(\operatorname{im}T\).
    \marginnote{To identify a span with the image, prove both inclusions.
      Then check linear independence to obtain a basis.}

  \item Both bases have two vectors, so
    \[
      \operatorname{nullity}T=2,
      \qquad
      \operatorname{rank}T=2.
    \]
    The rank--nullity identity is verified by
    \[
      \operatorname{rank}T+\operatorname{nullity}T
        =2+2=4=\dim\mathbb{R}^4.
    \]

  \item The map \(T\) is not injective, since its nullity is \(2\).
    It is not surjective, since
    \(\operatorname{rank}T=2<3=\dim\mathbb{R}^3\).
\end{enumerate}
\end{proof}
